\section{REFLECTION OF THE TEAM MEMBERS ON THE PROJECT}

Our work on the project ``Web Advertisement Optimization Using Reinforcement Learning'' was an insightful journey that gave us real-world exposure to applying reinforcement-learning ideas to a problem of clear commercial and societal importance. Throughout this project, we explored algorithm-driven strategies to optimise advertisement selection, gaining both technical knowledge and collaborative experience. Each of us took on specific responsibilities, which contributed to the overall success of the project and helped us grow in our respective areas. We held weekly synchronisation meetings, divided the work clearly across the algorithmic core, the environment, the front-end and the documentation, and reviewed each other's contributions before they were merged into the main branch. The result is a system that none of us could have built individually in the time available, and a report that reflects the integrated nature of our team-work.

\paragraph{Aditya Mishra (2241013120).}
I was responsible for the algorithmic core of the project. I designed the \texttt{BaseAgent} abstract class that defines the contract every concrete reinforcement-learning agent must satisfy, and I implemented all five concrete agents --- the Random baseline, $\varepsilon$-Greedy, Decaying $\varepsilon$-Greedy, UCB1 and Thompson Sampling --- against that interface. The most technically interesting parts were (i) implementing the Beta--Bernoulli posterior sampling for Thompson Sampling, which required writing a Gamma sampler using the Marsaglia--Tsang algorithm because no off-the-shelf statistics library was available in the browser environment, and (ii) carefully arranging the simulation loop so that every agent observes the same Bernoulli draw on every step. This work taught me the importance of clean abstractions and consistent interfaces when implementing several agents that all need to share the same simulation loop. I gained a deep understanding of the exploration--exploitation trade-off, the regret bounds of the various algorithms, and the practical engineering trade-offs between a deterministic policy (UCB1) and a stochastic one (Thompson Sampling). I also led the technical writing of Sections 3 and 4 of this report, including the pseudocode and the comparative analysis. The experience improved my TypeScript skills, my comfort with statistical estimation, and my ability to contribute efficiently to a multi-person engineering project.

\paragraph{Sarwan Samal (2241019260).}
My main responsibility was the environment design and the experimental harness. I implemented the \texttt{AdEnvironment} class, including the configurable Bernoulli ad-server, the CTR-mode and Revenue-mode reward signals, the optional sinusoidal drift for non-stationary experiments, and the oracle that computes the optimal expected reward for regret accounting. I also built the headless benchmark loop that lets us run the same scenario hundreds of times with different seeds without going through the React render pass, which was essential for producing the 100-run aggregate statistics reported in Section~4. Along the way, I tuned each agent's hyperparameters --- the exploration rate of $\varepsilon$-Greedy, the decay schedule of Decaying $\varepsilon$-Greedy, the confidence constant of UCB1, and the Beta prior of Thompson Sampling --- and learned how the choice of hyperparameter can change the qualitative behaviour of an agent. I learned how evaluation metrics like cumulative reward, cumulative regret and convergence step help verify performance, and how to design an experimental protocol that controls for the noise inherent in a stochastic environment. This hands-on work gave me a deeper understanding of statistical experimentation and of the difference between asymptotic theoretical guarantees and finite-time empirical performance.

\paragraph{Priyanka Sahoo (2241019109).}
I led the presentation layer of the project. I designed and implemented the four primary views of the application --- the Simulator, the Dashboard, the History view and the About / education page --- using React 19, Tailwind CSS and Radix UI. The Simulator view was the most demanding because it shows the agents' decisions in real time and must update at sixty frames per second without dropped frames; I learned how to use Zustand selectors and \texttt{React.memo} to keep the re-render cost minimal even when the simulation is producing many updates per second. I also built the four chart variants in Recharts --- the cumulative-reward chart, the cumulative-regret chart, the rolling-CTR chart and the exploration heat-map --- each tied directly to a slice of the live application state. Each chart had its own subtleties: the regret chart needed careful y-axis scaling so that the curves remained legible at small regret values, and the heat-map required a custom renderer because Recharts does not ship a built-in heat-map component. This role gave me a deep appreciation for how an effective visual representation can simplify complex information. It also improved my ability to interpret data and explain outcomes clearly to a non-specialist audience, which is an essential skill in analytics and in any data-science career.

\paragraph{Rishivatsal Mishra (2241013165).}
My contribution to this project was the literature survey, the references, the documentation and team coordination. I conducted the survey of the multi-armed bandit literature that is reported in Section~2, compiled the thirty academic references, and traced the genealogy of the five algorithms from Robbins (1952) and Thompson (1933) through Auer (2002) and Chapelle (2011) to recent contextual and deep-RL extensions. I authored Section~1 (Introduction), Section~5 (Conclusions and Future Scope) and this Reflection section, and I led the final pass of typesetting, cross-reference cleanup and proofreading. I also organised our weekly synchronisation meetings, ensured clear communication between the algorithmic, environmental and front-end work, and helped resolve issues during development. This role reinforced my ability to manage responsibilities inside a group, taught me how detailed documentation supports the transparency and reproducibility of a project, and gave me confidence in technical writing for an academic audience. The literature survey in particular was a substantial learning experience: reading and synthesising thirty papers required me to develop a working understanding of the entire bandit field, not just the algorithms our system implements.

\paragraph{Team-level reflections.}
Working together on this project taught us several lessons that none of us could have learnt alone. The first is that a clean, well-defined interface --- the \texttt{BaseAgent} abstract class --- enables parallel work in a way that ad-hoc co-ordination cannot. Once the interface was specified, the agent implementations, the environment, the simulation engine, and the chart components could all be developed independently and integrated with very little re-work. The second lesson is the importance of reproducibility: by persisting the random seed of every saved run we made our results auditable and re-runnable, which became valuable when reviewers asked us to verify specific numbers. The third lesson is that good visualisation is not a cosmetic concern but a debugging tool: several subtle bugs in the early agent implementations were spotted only when their internal state was rendered live in the dashboard. The fourth lesson is the value of small, focused commits: by keeping each Git commit limited to a single logical change, we made it possible to bisect bugs to the exact line of code that introduced them, which paid off twice during the project --- once when a refactor of the simulation engine introduced a subtle off-by-one error in the regret accounting, and once when a CSS change inadvertently broke the responsive layout of the History view.

\paragraph{Skills gained during the project.}
Beyond the project-specific outcomes, each member acquired skills that will be applicable far beyond the bandit-simulation domain:
\begin{itemize}\itemsep0pt
  \item \textbf{Statistical thinking.} We learned to distinguish asymptotic from finite-time guarantees, to think about variance as well as mean performance, to design experiments that control for noise, and to read and apply formal regret bounds. The Wilcoxon and bootstrap intervals reported in Section~4.3.13 came directly from a course module on inferential statistics that we previously thought was disconnected from software engineering.
  \item \textbf{Software architecture.} We learned the practical value of layered architecture, dependency inversion, and explicit extension hooks. The fact that a sixth agent could be added in roughly thirty lines of code without touching the UI is the result of design decisions made on day one, not week ten.
  \item \textbf{Type-driven development.} We learned to express invariants in TypeScript types (discriminated unions, exhaustive switch statements, optional fields) and to lean on the compiler to enforce them. Several classes of runtime bug --- forgotten reward modes, missing arm indices, mis-typed callback shapes --- were eliminated entirely by the type system before we even reached the testing phase.
  \item \textbf{Technical writing.} We learned how to structure a long-form report, how to write pseudocode that reads as cleanly as code, how to design tables that communicate at a glance, and how to use cross-references and labels to keep a 70-page document navigable. The LaTeX skills acquired during this project will carry over to dissertations, conference papers and any future technical-writing tasks.
  \item \textbf{Project management.} We learned to plan in three-week sprints, to triage scope-creep, to maintain a backlog, and to track our own velocity against an agreed-upon timeline. The "what we did not do and why" section of this report (Section~5.2) is, in our view, as important as the "what we did" section --- and recognising that is itself a learned skill.
\end{itemize}

\paragraph{What we would do differently.}
With hindsight, three decisions would have benefited from earlier reconsideration. First, we should have written the headless benchmark harness on day one rather than week six; once it existed, the iteration speed of the entire experimental campaign at least doubled. Second, we should have committed to the C7-style report format earlier --- we initially drafted an appendix-heavy first version and had to restructure midway through, which cost roughly a week of writing time. Third, we should have set up a continuous-integration job to compile the LaTeX on every push; manually compiling locally meant that small markup errors sometimes slipped in unnoticed for several commits.

\paragraph{Looking forward.}
The five algorithms implemented in this project are a small subset of the modern bandit-and-RL landscape. Each of us is leaving the project with a clear personal next-step: Aditya plans to extend the agent pool with a contextual variant (LinUCB) for his Master's research, Sarwan intends to write a follow-up note on the reward-mismatch finding for an undergraduate workshop, Priyanka is exploring how the same React + Zustand pattern could be applied to other interactive scientific visualisations, and Rishivatsal will be using the LaTeX scaffolding from this report as a template for future technical writing. The simulator itself will remain hosted online for at least the duration of the academic year, so that future students at the institute can use it as a teaching aid for the multi-armed bandit problem.

In summary, this project helped us observe how reinforcement-learning algorithms behave in a realistic context while developing valuable skills in teamwork, problem-solving and technical communication. The experience has prepared us well for future academic and professional challenges in the field of computer science and engineering, and we are grateful to our supervisor, our coordinators, and our peers for the support that made this work possible.
\newpage
